\section{Introducción}

En este informe se presentan dos experiencias experimentales orientadas a determinar magnitudes físicas de gran relevancia: la densidad de un cuerpo de forma cilíndrica y la aceleración de la gravedad . A lo largo de ambos procedimientos se trabajó con mediciones directas e indirectas, considerando las incertidumbres asociadas y analizando de qué manera estas afectan los resultados obtenidos. Asimismo, el desarrollo de las experiencias permitió poner en práctica herramientas de análisis experimental, interpretación de datos y estudio de relaciones entre variables físicas.

En la primera parte del trabajo se llevó a cabo la determinación la densidad de un cilindro a partir de la medición de su masa y de sus dimensiones geométricas. La altura y el diámetro fueron medidos utilizando instrumentos de diferente precisión, específicamente una regla graduada y un calibre Vernier, mientras que la masa se obtuvo con una balanza de mayor exactitud. Esta comparación permitió evaluar la influencia de la precisión instrumental en el cálculo del volumen y, en consecuencia, en la determinación final de la densidad, recurriendo para ello a criterios de estimación y propagación de incertidumbres.

En el segundo ejercicio se buscó determinar de la aceleración de la gravedad mediante el uso de un péndulo simple. Para ello, se registró el período de oscilación correspondiente a distintas longitudes de hilo y, a partir de esos datos, se estimó el valor de g. Durante el análisis se tuvieron en cuenta las incertidumbres de las mediciones de tiempo y longitud, así como también las limitaciones propias del modelo teórico empleado para describir el movimiento del péndulo en condiciones reales de laboratorio.

En síntesis, ambos experimentos no solo permitieron obtener valores experimentales de la aceleración de la gravedad y de la densidad del cilindro, sino que también sirvieron para profundizar en el tratamiento de errores, el alcance de los modelos teóricos y el análisis cuantitativo de fenómenos físicos observados en el laboratorio.